% Folder 33, batch 08 — archivists' pages 141–160.
% Pass: Opus 5.5 (claude-opus-5-5), 2026-10-03 — first pass, unchecked against the pages by a human.
%
% Only the odd pages carry his mathematics; the even pages are versos of
% unrelated typescripts and are not transcribed (the gaps in the numbering are
% the record). Notation fixed for the batch: complexes written F. are set
% $F_\bullet$; his double-struck R with script Hom is set $\mathbb{R}\mathcal{H}om$;
% his « Hom° » (degree 0) is set $\mathrm{Hom}^{\circ}$.
\documentclass[11pt,a4paper]{article}
\input{../preamble/grothendieck}

\folder{33}
\batch{8}
\pages{141}{160}
\foldertitle{SGA 1965. MS [Manuscrit] brouillon : notes manuscrites (s.d.).}
\dating{1965-[vers 1971]}
\watermark{Édition de démonstration}

\begin{document}

\page{141}
\note{la page commence au milieu d'un argument venu des pages précédentes (majoration de l'amplitude cohomologique de $R_!f$), qui ne figure pas dans ce lot.}

$R_!f(L_\bullet)$ est de tor-dim. majorée par $[N-2d, M]$ \struck{\ill{}}. (utilisant \ill{} \ill{} déjà \uncertain{finitude} !)

Résultat formel de Künneth : \ill{}.
\[
R_!f(L_\bullet)\overset{L}{\otimes} L'_\bullet \;\simeq\; R_!f\bigl(L_\bullet \overset{L}{\otimes} Lf^*(L'_\bullet)\bigr) \qquad 2.
\]
\note{entre $\overset{L}{\otimes}$ et $L'_\bullet$, un symbole biffé.}

$L'_\bullet$ est de \add{\uncertain{tor-}}dim. majorée par $[N', M']$, $Lf^*(L'_\bullet)$ itou, $L_\bullet\overset{L}{\otimes} Lf^*(L'_\bullet)$ est à coh. de degrés majorés par $[N+N', M+M']$, donc $R_!f(\;\;)$ est à coh. de degrés bornés par (\uncertain{condition} \uncertain{homologique} !) $[N+N'-2d, M+M']$, ce qui est la relation qu'on voulait prouver.

\struck{\ill{}} \note{suivent trois lignes biffées d'un trait et d'une boucle, en partie illisibles ; on y lit encore \uncertain{supposons}, \uncertain{conditions pour simplifier}, et, dans un cadre, \uncertain{Conditions}.}

\page{143}
\section{(32) Formule de Lefschetz-Verdier}
\note{le titre et le numéro cerclé 32 sont de sa main, en tête de page.}

\begin{tikzcd}[column sep=small, row sep=small]
 & X\times_S Y \arrow[dl, "\mathrm{pr}_1"'] \arrow[dr, "\mathrm{pr}_2"] \arrow[dd, "h"] & \\
F_\bullet\;\; X \arrow[dr, "f"'] & & Y\;\; G_\bullet \arrow[dl, "g"] \\
 & S &
\end{tikzcd}

(a)
\[
\struck{\text{Ext}}\;\mathrm{Hom}^{\circ}\bigl(\mathrm{pr}_1^*(F_\bullet), \mathrm{pr}_2^!(G_\bullet)\bigr) \simeq R\Gamma_{X\times_S Y}\,\mathbb{R}\mathcal{H}om^{\bullet}\bigl(\mathrm{pr}_1^*(F_\bullet), \mathrm{pr}_2^!(G_\bullet)\bigr)
\]
\[
\begin{aligned}
&\wr\;\; \mathrm{Hom}^{\circ}\bigl(\mathrm{pr}_{2!}\,\mathrm{pr}_1^*(F_\bullet), G_\bullet\bigr)\\
&\wr\;\; \mathrm{Hom}^{\circ}\bigl(g^*f_!(F_\bullet), G_\bullet\bigr)\\
&\wr\;\; \mathrm{Hom}^{\circ}\bigl(f_!(F_\bullet), g_*(G_\bullet)\bigr)
\end{aligned}
\]
\marginal{correspondances cohomologiques de degré 0 de $(X,F_\bullet)$ dans $(Y,G_\bullet)$ sur $S$ …}

\[
\begin{aligned}
\mathrm{Hom}^{\circ}\bigl(\mathrm{pr}_1^*(F_\bullet), \mathrm{pr}_2^!(G_\bullet)\bigr) &\simeq \mathrm{Hom}^{\circ}\bigl(f_!(F_\bullet), g_*(G_\bullet)\bigr)\\
\mathrm{Hom}^{\circ}\bigl(\mathrm{pr}_2^*(G_\bullet), \mathrm{pr}_1^!(F_\bullet)\bigr) &\simeq \mathrm{Hom}^{\circ}\bigl(g_!(G_\bullet), f_*(F_\bullet)\bigr)
\end{aligned}
\]

Supposons que $f$ et $g$ soient propres, alors les deux expressions précédentes s'accouplent dans \struck{$\mathrm{Hom}^{\circ}\,Rf_*$} $\mathrm{Hom}^{\circ}(f_*(F_\bullet), f_!(F_\bullet))$ et dans $\mathrm{Hom}^{\circ}(g_*(G_\bullet), g_!(G_\bullet))$ (si $S$ discret, \add{et $F_\bullet$ et $G_\bullet$ \uncertain{de tor-dim. finie}}) \struck{\ill{}}, donc \struck{\ill{}} dans $A_S$ (par « \uncertain{nombres} \uncertain{de} \uncertain{Lefschetz} » ou « \uncertain{traces} »).

\note{les lignes qui suivent sont barrées de longs traits obliques ; elles restent lisibles en partie.}

\struck{Cet accouplement est induit} \struck{d'un accouplement de nature locale} \struck{(par le fait} \struck{\uncertain{que}} \struck{$f$ et $g$ \uncertain{soient} propres)}

\[
\text{\struck{par}}\quad \mathbb{R}\mathcal{H}om^{\bullet}\bigl(\mathrm{pr}_1^*(F_\bullet), \mathrm{pr}_2^!(G_\bullet)\bigr) \times \mathbb{R}\mathcal{H}om\bigl(\mathrm{pr}_2^*(G_\bullet), \mathrm{pr}_1^!(F_\bullet)\bigr) \longrightarrow h^!(A_S)
\]
qu'on définit ainsi \ill{}

\note{le passage depuis « Cet accouplement » est encadré, et l'encadré est en partie biffé ; la définition annoncée est reprise p.~147.}

\page{145}
(b) $\mathcal{L}_\bullet$ un complexe pseudo-cohérent de Tor-dimension finie. Alors

\begin{tikzcd}
A_X \arrow[r, "\varepsilon"] \arrow[drr, "\tau_{\mathcal{L}_\bullet}"'] & \mathbb{R}\mathcal{H}om(\mathcal{L}_\bullet, \mathcal{L}_\bullet) & \check{\mathcal{L}}_\bullet \overset{L}{\otimes} \mathcal{L}_\bullet \arrow[l, "\simeq"'] \arrow[d, "\tau"] \\
 & & A_X
\end{tikzcd}

\note{l'indice de $\tau_{\mathcal{L}_\bullet}$ surcharge un premier indice biffé.}

Cet accouplement induit un homomorphisme \uncertain{\ill{}}
\[
\mathcal{H}^{0}\bigl(\mathbb{R}\mathcal{H}om(\mathcal{L}_\bullet, \mathcal{L}_\bullet)\bigr) \longrightarrow A_X
\]
qui s'explicite, lorsque $\mathcal{L}_\bullet$ est \struck{plat} \add{loc. libre de type fini} \struck{\ill{}} en chaque degré, comme
\[
\varphi \longmapsto \sum (-1)^i\, \mathrm{Tr}\, \varphi_i
\]
où $\varphi = (\varphi_i)$ est \struck{\ill{}} un \uncertain{homomorphisme} de complexes $\mathcal{L}_\bullet \to \mathcal{L}_\bullet$. [\struck{Regardons en effet}

\struck{$\varphi_i \in \check{\mathcal{L}}_i\otimes\mathcal{L}_i$, \uncertain{on} \uncertain{doit} \uncertain{les} \uncertain{accoupler}}

Supposons maintenant que $\mathcal{L}_\bullet$, $\mathcal{M}_\bullet$ \add{\uncertain{soient}} de Tor-dimension finie et pseudo-cohérents, considérons les \uncertain{couplages} évidents
\[
\begin{aligned}
\mathbb{R}\mathcal{H}om(\mathcal{L}_\bullet, \mathcal{M}_\bullet) \times \mathbb{R}\mathcal{H}om(\mathcal{M}_\bullet, \mathcal{L}_\bullet) &\xrightarrow{\;\varphi_{\mathcal{L}_\bullet,\mathcal{M}_\bullet}\;} \mathbb{R}\mathcal{H}om(\mathcal{L}_\bullet, \mathcal{L}_\bullet) & (u,v)&\mapsto vu\\
\mathbb{R}\mathcal{H}om(\mathcal{M}_\bullet, \mathcal{L}_\bullet) \times \mathbb{R}\mathcal{H}om(\mathcal{L}_\bullet, \mathcal{L}_\bullet) &\longrightarrow \mathbb{R}\mathcal{H}om(\mathcal{M}_\bullet, \mathcal{M}_\bullet) & (u,v)&\mapsto uv
\end{aligned}
\]
\note{au second facteur de la deuxième ligne, la page porte bien $(\mathcal{L}_\bullet, \mathcal{L}_\bullet)$, là où l'on attend $(\mathcal{L}_\bullet, \mathcal{M}_\bullet)$.}

alors \struck{les \ill{}} on a
\[
\tau_{\mathcal{L}_\bullet}\,\varphi_{\mathcal{L}_\bullet,\mathcal{M}_\bullet} = \tau_{\mathcal{M}_\bullet}\,\varphi_{\mathcal{M}_\bullet,\mathcal{L}_\bullet}^{\,\mathrm{sym}}
\]

\page{147}
et ses accouplements \struck{induisent des} \struck{\ill{}}
\[
\mathbb{R}\mathcal{H}om(\mathcal{L}_\bullet, \mathcal{M}_\bullet) \times \mathbb{R}\mathcal{H}om(\mathcal{M}_\bullet, \mathcal{L}_\bullet) \longrightarrow A_X
\]
identifient chacun des deux \uncertain{faisceaux} $P_\bullet$, $Q_\bullet$ au \uncertain{premier} \uncertain{membre} au « dual » de l'autre
\[
P_\bullet \simeq \mathbb{R}\mathcal{H}om(Q_\bullet, A_X), \qquad Q_\bullet \simeq \mathbb{R}\mathcal{H}om(P_\bullet, A_X^\bullet)
\]

Du point de vue \uncertain{valeurs}, si $u^{(i)} : \mathcal{L}_\bullet \to \mathcal{M}_\bullet[i]$, $v^{(-i)} : \mathcal{M}_\bullet \to \mathcal{L}_\bullet[-i]$, donc \struck{\ill{}} $v^{(-i)}u^{(i)} : \mathcal{L}_\bullet \to \mathcal{M}_\bullet$.
\note{sic : on attend $\mathcal{L}_\bullet \to \mathcal{L}_\bullet$.}
et
\[
\langle u^{(i)}, v^{(-i)}\rangle = \sum \mathrm{Tr}\,(uv)^{\varepsilon}_i\,(-1)^i
\]

(c) Dans le cas où, \add{$f$ et $g$ propres}, supposons que $S$ est \struck{\uncertain{de dimension finie}} discret, $F_\bullet$, $G_\bullet$ à coh. constructible \add{à Tor-dim. finie} \struck{\ill{}}, \add{donc \ill{} \uncertain{finie}}, \uncertain{j'ai dit}
que $\mathrm{Hom}^{\circ}(Rf_*(F_\bullet), Rg_*(G_\bullet))$ et $\mathrm{Hom}^{\circ}(Rg_*(G_\bullet), Rf_*(F_\bullet))$ \add{ou $\mathbb{R}\mathrm{Hom}(-,\cdot)\times\mathbb{R}\mathrm{Hom}(\cdot,-)$ \ill{}} sont \uncertain{duals} l'un de l'autre, \uncertain{et} \uncertain{je} \uncertain{dis} que cette dualité est \uncertain{induite} par un

accouplement
\[
\begin{aligned}
&\mathbb{R}\mathcal{H}om\bigl(\mathrm{pr}_1^*(F_\bullet), \mathrm{pr}_2^!(G_\bullet)\bigr) \times \mathbb{R}\mathcal{H}om\bigl(\mathrm{pr}_2^*(G_\bullet), \mathrm{pr}_1^!(F_\bullet)\bigr)\\
&\qquad\longrightarrow K^\bullet_{X\times_S Y} = h^!(A_S)
\end{aligned}
\]
qu'on définit ainsi (sans hyp. de propreté sur $f$, $g$ !) On a, si $Z = X\times_S Y$

\page{149}
\[
\mathrm{pr}_2^!(G_\bullet) \longrightarrow D_Z\,\mathrm{pr}_2^*(D_Y G_\bullet)
\]
\note{au-dessus de $\mathrm{pr}_2^!(G_\bullet)$, une accolade porte $D_Y(D_Y(G_\bullet))$ ; sous la flèche : « iso si $D_Y(G_\bullet)$ à coh. constr. ». Le $G_\bullet$ surcharge un $F_\bullet$ biffé ; devant $D_Z$, un symbole biffé illisible.}

d'où
\[
\begin{aligned}
\mathbb{R}\mathcal{H}om\bigl(\mathrm{pr}_1^*(F_\bullet), \mathrm{pr}_2^!(G_\bullet)\bigr) &\longrightarrow \mathbb{R}\mathcal{H}om\bigl(\mathrm{pr}_1^*(F_\bullet), D_Z\,\mathrm{pr}_2^*(D_Y G_\bullet)\bigr)\\
&\;\wr\;\; D_Z\bigl(\mathrm{pr}_1^*(F_\bullet)\overset{L}{\otimes}\mathrm{pr}_2^*(D_Y G_\bullet)\bigr)
\end{aligned}
\]

Il suffit de décrire maintenant
\[
D_Z\bigl(\mathrm{pr}_1^*(F_\bullet)\overset{L}{\otimes}\mathrm{pr}_2^*(D_Y(G_\bullet))\bigr) \longrightarrow D_Z\,\mathbb{R}\mathcal{H}om\bigl(\mathrm{pr}_2^*G_\bullet, \mathrm{pr}_1^!(F_\bullet)\bigr)
\]
i.e. un \uncertain{morphisme} \ill{}
\[
\mathbb{R}\mathcal{H}om\bigl(\mathrm{pr}_2^*(G_\bullet), \mathrm{pr}_1^!(F_\bullet)\bigr) \longrightarrow \mathrm{pr}_1^*(F_\bullet)\overset{L}{\otimes}\mathrm{pr}_2^*(D_Y(G_\bullet))
\]
Or le premier membre s'envoie \ill{}, comme \ill{} \ill{}, dans
\[
D_Z\bigl(\mathrm{pr}_2^*(G_\bullet)\overset{L}{\otimes}\mathrm{pr}_1^*(D_X F_\bullet)\bigr)
\]
\ill{} il faut \uncertain{définir}
\[
D_Z\bigl(\mathrm{pr}_2^*(G_\bullet)\overset{L}{\otimes}\mathrm{pr}_1^* D_X(F_\bullet)\bigr) \longrightarrow \mathrm{pr}_1^*(F_\bullet)\overset{L}{\otimes}\mathrm{pr}_2^*(D_Y(G_\bullet))
\]
\note{sous cette flèche, une flèche courbe en retour, de droite à gauche.}
\ill{}, \uncertain{définissons} \ill{} \ill{}.

Pour ceci, \ill{}, \ill{} \uncertain{et} \ill{} \ill{}

\ill{} \uncertain{isomorphisme} \ill{} \ill{} ! \note{les deux lignes précédentes sont soulignées d'un double trait.}

On \struck{\ill{}} \add{\ill{}} \ill{} \ill{} \uncertain{isomorphisme} \ill{} \uncertain{définition}

\ill{} \ill{} \ill{}

\[
(\ast)\quad \mathrm{pr}_1^*(F_\bullet)\overset{L}{\otimes}\mathrm{pr}_2^* D_Y G_\bullet \overset{L}{\otimes} \mathrm{pr}_2^* G_\bullet \overset{L}{\otimes} \mathrm{pr}_1^* D_X(F_\bullet) \longrightarrow K^\bullet_Z
\]
\note{le signe $(\ast)$ est cerclé ; il renvoie à la marge. Devant $K^\bullet_Z$, un symbole biffé illisible.}

\marginal{Il faut \uncertain{voir} \uncertain{que} \ill{} \uncertain{deux} \ill{} \uncertain{dans} $K^\bullet_Z$ \ill{}
\[
\begin{aligned}
\mathbb{R}\mathcal{H}om(\mathrm{pr}_1^*F_\bullet, \mathrm{pr}_2^!G_\bullet) &\to D_Z\bigl(\mathrm{pr}_1^*(F_\bullet)\overset{L}{\otimes}\mathrm{pr}_2^*(D_Y G_\bullet)\bigr)\\
\mathbb{R}\mathcal{H}om(\mathrm{pr}_2^*G_\bullet, \mathrm{pr}_1^!F_\bullet) &\to D_Z\bigl(\mathrm{pr}_2^*(G_\bullet)\overset{L}{\otimes}\mathrm{pr}_1^*(D_X F_\bullet)\bigr)
\end{aligned}
\]
\uncertain{Plus} \ill{}. (Note écrite le long du bord gauche, avec un signe $(\ast)$ cerclé.)}

\page{151}
On \uncertain{a} \ill{}
\[
\begin{aligned}
\mathrm{pr}_1^*(F_\bullet)\otimes\mathrm{pr}_1^*(D_X F_\bullet) &\longrightarrow \mathrm{pr}_1^*(R^\bullet_X)\\
\mathrm{pr}_2^*(G_\bullet)\otimes\mathrm{pr}_2^* D_Y G_\bullet &\longrightarrow \mathrm{pr}_2^*(R^\bullet_Y)
\end{aligned}
\]
\uncertain{d'où} \uncertain{le} \uncertain{premier} \uncertain{membre} de $(\ast)$. \ill{} \ill{}.

\ill{} dans
\[
\mathrm{pr}_1^*(R^\bullet_X)\otimes\mathrm{pr}_2^*(R^\bullet_Y)
\]
il faut l'envoyer \add{\uncertain{canoniquement}} (\ill{} dans) $R^\bullet_Z$. Ceci \uncertain{devrait} \uncertain{par} \uncertain{effet} \ill{} \ill{} \ill{} \ill{}, \uncertain{p. ex.} \uncertain{s'il} \ill{} \uncertain{trivial} \uncertain{si} $X$ ou $Y$ \uncertain{lisse}$/k$.

[\ill{} \ill{}, \ill{} \ill{}. Que \ill{} \ill{} \ill{} $\to \sim K^\bullet_Z$ ]. Cet \ill{} \uncertain{définit} bien l'accouplement \uncertain{voulu}, d'où un \ill{}

\marginal{\uncertain{Même} si $X$ ou $Y$ \uncertain{lisse} sur $k$, \ill{} \ill{} \ill{} \uncertain{des} \uncertain{résolutions} \uncertain{des} \uncertain{singularités} \ill{} !}

\uncertain{Formule} \uncertain{cruciale}.
\[
(\ast\ast)\quad \mathrm{pr}_1^*(F_\bullet)\otimes\mathrm{pr}_2^*(D_Y G_\bullet) \longrightarrow D_Z\bigl(\mathrm{pr}_2^* G_\bullet\otimes\mathrm{pr}_1^* D_X(F_\bullet)\bigr)
\]
\note{le signe $(\ast\ast)$ est cerclé, et surcharge peut-être autre chose.}

Le théorème fondamental de Verdier : \uncertain{Cet} \uncertain{homomorphisme} \uncertain{est} un isomorphisme.

\drawing{au bas de la page, un réseau de lignes obliques entrecroisées (une rangée de losanges) sur lequel il a écrit en travers : « Dire », « $X$, $Y$ \uncertain{lisses} », « $F_\bullet$, $G_\bullet$ », « \uncertain{revient} \uncertain{à} \uncertain{dire} … \uncertain{d'habitude} », « $i_!(A_U)$, $j_!(A_V)$ », « $i : U\to X$, $j : V\to X$ », « $U$, $V$ », « \ill{} ». Le dessin et l'écriture paraissent superposés ; l'ordre de lecture n'est pas sûr.}

\page{153}
L'existence d'un \uncertain{homom.} dans ce cas \uncertain{général} \uncertain{donc}.
\[
(1)\quad R'_X \overset{L}{\boxtimes} R'_Y \xrightarrow{\;\sim\;} R'_{X\times Y}
\]
et du fait que l'accouplement de Verdier \uncertain{donne} un isomorphisme
\[
(2)\quad F_\bullet \overset{L}{\boxtimes} D_Y G_\bullet \xrightarrow{\;\sim\;} D_Z\bigl(D_X F_\bullet \overset{L}{\boxtimes} G_\bullet\bigr)
\]
\note{dans (2), un symbole biffé devant $F_\bullet$ et un autre devant $D_X F_\bullet$.}
résulte \uncertain{de} \uncertain{la} \uncertain{formule} générale \uncertain{cohomologique} \uncertain{suivante} :

\marginal{« Künneth local »}

\[
\mathbb{R}\mathcal{H}om\bigl(F_\bullet \overset{L}{\boxtimes} G_\bullet,\, P_\bullet \overset{L}{\boxtimes} Q_\bullet\bigr) \xleftarrow{\;\simeq\;} \mathbb{R}\mathcal{H}om(F_\bullet, P_\bullet) \overset{L}{\boxtimes} \mathbb{R}\mathcal{H}om(G_\bullet, Q_\bullet)
\]
\[
R(f\times g)^!\bigl(P_\bullet \overset{L}{\boxtimes} Q_\bullet\bigr) \simeq Rf^!(P_\bullet) \overset{L}{\boxtimes} Rg^!(Q_\bullet)
\]
\note{la première formule est encadrée ; deux traits obliques la relient à la seconde. Dans la seconde, la place des signes $!$ n'est pas sûre.}

valable si $F_\bullet$, $G_\bullet$ sont \uncertain{de} \uncertain{Tor-dim.} \uncertain{finie} \add{$\in \mathrm{Ob}\,D^{-}$} \struck{\ill{}} à coh. constructible, et $P_\bullet$, $Q_\bullet$ \struck{\uncertain{tor-dim.}} ($A = \mathbb{Z}/n\mathbb{Z}$, $n$ premier à car. $k$),
\struck{\uncertain{finie}} \uncertain{faudrait} valable si \uncertain{on} \uncertain{passe} \uncertain{par} les schémas $U$, $V$ (la résolution des singularités, et aux \uncertain{précédents} \uncertain{sur} eux, \struck{\ill{}}
\uncertain{étant} possible : O.K. \uncertain{en} car. 0, \uncertain{et} \uncertain{en} car. $p$ \uncertain{quelc.} si $U$, $V$ de dim.
$\leq 2$). N.B. Dans la \uncertain{démonstration} de (1), \uncertain{on} \uncertain{se} \uncertain{ramène} au cas \uncertain{où} $U$, $V$ \uncertain{des} schémas lisses
qui contiennent $X$, $Y$ — \uncertain{ce} qui introduit des

\page{155}
dimensions \uncertain{parasites}, qui peuvent \uncertain{compliquer} les \uncertain{questions} de \uncertain{résolution}. Il serait préférable \struck{\ill{}} (si $X$, $Y$ ne sont pas \uncertain{eux-mêmes} lisses) de \struck{\ill{}} \add{pouvoir} utiliser localement des \add{\uncertain{finis}} \uncertain{homomorphismes} $X\to X'$, $Y\to Y'$, avec $X'$ ($Y'$) lisse de dimension \struck{\ill{}} \uncertain{égale} \uncertain{à} $\dim X$ ($\leq \dim Y$)
et des \uncertain{morphismes} \uncertain{qui} \uncertain{garantissent} \uncertain{l'interprétation} de $\varphi^!$, $\psi^!$ \struck{des} \uncertain{par} \uncertain{des} \uncertain{morphismes} finis $\varphi$, $\psi$,

\marginal{en marge gauche :
$X \xrightarrow{\;\varphi\;} X'$, \quad $Y \xrightarrow{\;\psi\;} Y'$ (flèches verticales).}

on a bien
\[
(\varphi\times\psi)^!\;\text{\struck{$\varphi^!$}}\;\bigl(P_\bullet \boxtimes Q_\bullet\bigr) \simeq \varphi^!(P_\bullet)\boxtimes\psi^!(Q_\bullet)
\]
\uncertain{ce} qui \uncertain{montre}
\[
\mathbb{R}\mathcal{H}om\bigl(F_\bullet\boxtimes G_\bullet,\, P_\bullet\boxtimes Q_\bullet\bigr) \simeq \mathbb{R}\mathcal{H}om(F_\bullet, P_\bullet)\overset{L}{\boxtimes}\mathbb{R}\mathcal{H}om(G_\bullet, Q_\bullet)
\]
\uncertain{pour} $F_\bullet = \varphi_*(A_X)$, $G_\bullet = \psi_*(A_Y)$ ----

\page{157}
(d) \struck{Cela posé}

On a donc trouvé un accouplement (\uncertain{dorénavant} : la \uncertain{restriction} \ill{} \ill{} \ill{})
\[
\mathbb{R}\mathcal{H}om\bigl(\mathrm{pr}_1^*(F_\bullet), \mathrm{pr}_2^!(G_\bullet)\bigr) \times \mathbb{R}\mathcal{H}om\bigl(\mathrm{pr}_2^*(G_\bullet), \mathrm{pr}_1^!(F_\bullet)\bigr) \longrightarrow K^\bullet_{X\times_S Y}
\]
\struck{D'où, pour} \struck{$\xi^\circ \in \mathbb{R}^0\Gamma_{X\times_S Y}(\ldots)$} \struck{et} \struck{$\eta^\circ \in \mathbb{R}^0\Gamma_{X\times_S Y}(\ldots)$}
\struck{un élément} \struck{$\xi^\circ\cup\eta^\circ \in \mathbb{R}^0\Gamma_{X\times_S Y}(K^\bullet_{X\times_S Y})$}
\note{dans la ligne biffée, deux traits courbes renvoient les parenthèses vides aux deux facteurs de l'accouplement ci-dessus.}

Appliquant $\mathbb{R}\Gamma_{X\times_S Y}$, \uncertain{on} \uncertain{trouve} \uncertain{un} accouplement

\begin{tikzcd}[column sep=small, row sep=small, nodes={font=\scriptsize}]
\mathbb{R}\mathrm{Hom}^{\bullet}\bigl(\mathrm{pr}_1^*(F_\bullet), \mathrm{pr}_2^!(G_\bullet)\bigr) \times \mathbb{R}\mathrm{Hom}^{\bullet}\bigl(\mathrm{pr}_2^*(G_\bullet), \mathrm{pr}_1^!(F_\bullet)\bigr) \arrow[r] \arrow[d, no head, "\wr"] & \mathbb{R}K^\bullet_{X\times_S Y} \arrow[d] \\
\mathbb{R}\mathrm{Hom}^{\bullet}\bigl(f_*(F_\bullet), g_*(G_\bullet)\bigr) \times \mathbb{R}\mathrm{Hom}^{\bullet}\bigl(g_*(G_\bullet), f_*(F_\bullet)\bigr) \arrow[r] & A_S
\end{tikzcd}

\note{à la deuxième ligne, $f_*$ et $g_*$ surchargent des $\mathrm{pr}$ biffés ; à droite, le $\mathbb{R}$ devant $K^\bullet$ est peut-être biffé. Le signe $\wr$ est écrit au-dessus de chacun des deux facteurs.}

qui ne doit être autre que l'accouplement déjà connu \uncertain{relatif} à $P_\bullet \simeq \mathrm{pr}_{1*}(F_\bullet)$ et $Q_\bullet \simeq \mathrm{pr}_{2*}(G_\bullet)$.

En particulier, si
\[
\begin{aligned}
\xi^\circ &\in \mathbb{R}^0\mathrm{Hom}\bigl(\mathrm{pr}_1^*(F_\bullet), \mathrm{pr}_2^!(G_\bullet)\bigr) = \mathbb{R}^0\Gamma_{X\times_S Y}\,\mathbb{R}\mathcal{H}om\bigl(\mathrm{pr}_1^*(F_\bullet), \mathrm{pr}_2^!(G_\bullet)\bigr)\\
\eta^\circ &\in \mathbb{R}^0\mathrm{Hom}\bigl(\mathrm{pr}_2^*(G_\bullet), \mathrm{pr}_1^!(F_\bullet)\bigr) = \mathbb{R}^0\Gamma_{X\times_S Y}\,\mathbb{R}\mathcal{H}om\bigl(\mathrm{pr}_2^*(G_\bullet), \mathrm{pr}_1^!(F_\bullet)\bigr)
\end{aligned}
\]
alors
\[
\xi^\circ\cup\eta^\circ \in \mathbb{R}^0\Gamma_{X\times_S Y}\bigl(K^\bullet_{X\times_S Y}\bigr) \xrightarrow{\;\mathrm{Tr}_h\;} A.
\]
et on aura
\[
\boxed{\;\mathrm{Tr}(\xi^\circ\cup\eta^\circ) = \langle \xi^\circ, \eta^\circ\rangle\;}
\qquad \Bigl(= \sum_i \mathrm{Tr}\bigl((\xi^\circ)^{(i)}(\eta^\circ)^{(-i)}\bigr)(-1)^i\Bigr)
\]
\marginal{si on représente \struck{\ill{}} $f(F_\bullet)$, $g(G_\bullet)$ par des complexes plats de type fini en chaque degré.}

\page{159}
En particulier, si $X = Y$, $F_\bullet = G_\bullet$, prenant $\eta^\circ = $ \struck{\ill{}} $\delta_{X,F}$
\uncertain{on} \uncertain{trouve}
\[
\boxed{\;\mathrm{Tr}\bigl(\xi^\circ\cup\delta^\bullet_{X,F}\bigr) = \sum (-1)^i\,\mathrm{Tr}\,[\xi^\circ]^i\;}
\]
\note{sous la formule encadrée, deux petites marques indistinctes.}

C'est la formule de Lefschetz-Verdier.

(e) \struck{\ill{}} \note{sous le (e) cerclé, un petit « 2 ».}

Considérons un morphisme
\[
u : T \longrightarrow X\times_S Y
\]
explicité par deux morphismes
\[
u_1 : T\to X, \qquad u_2 : \text{\struck{$T$}}\;T\to Y
\]
On a alors un morphisme évident
\[
\mathbb{R}u_*\bigl(\mathbb{R}\mathcal{H}om(u_1^*(F_\bullet), u_2^!(G_\bullet))\bigr) \longrightarrow \mathbb{R}\mathcal{H}om\bigl(\mathrm{pr}_1^*(F_\bullet), \mathrm{pr}_2^!(G_\bullet)\bigr)
\]
Car $u_i = \mathrm{pr}_i u$, donc le premier membre \uncertain{est}

\begin{tikzcd}[column sep=small, row sep=small, nodes={font=\scriptsize}]
\mathbb{R}u_*\bigl(\mathbb{R}\mathcal{H}om(u^*\mathrm{pr}_1^*(F_\bullet), u^!(\mathrm{pr}_2^!(G_\bullet)))\bigr) \arrow[r, "\sim"] \arrow[d, no head, "\wr"] & \mathbb{R}\mathcal{H}om\bigl(\mathrm{pr}_1^*(F_\bullet), u_*u^!(\mathrm{pr}_2^!(G_\bullet))\bigr) \arrow[d, "\varepsilon"] \\
\mathbb{R}\mathcal{H}om\bigl(u_*u^*\mathrm{pr}_1^*(F_\bullet), \mathrm{pr}_2^!(G_\bullet)\bigr) \arrow[r] & \mathbb{R}\mathcal{H}om\bigl(\mathrm{pr}_1^*(F_\bullet), \mathrm{pr}_2^!(G_\bullet)\bigr)
\end{tikzcd}

\note{dans le premier terme, $u^*$ surcharge un $\mathrm{pr}$ biffé. Le signe vertical de la colonne de gauche est un $\wr$ incliné ; on le donne tel quel.}

Donc tout élément \struck{$\mathrm{Corr}^\circ(X,F)$}
\[
\alpha^\circ \in \mathbb{R}\mathrm{Hom}^\circ\bigl(u_1^*(F_\bullet), u_2^!(G_\bullet)\bigr) = \mathbb{R}^0\Gamma_T\bigl(\mathbb{R}\mathcal{H}om(u_1^*(F_\bullet), u_2^!(G_\bullet))\bigr)
\]
définit un élément
\[
\overline{\alpha^\circ} \in \mathrm{Corr}(X,F;\, Y,G)
\]
De même, un élément $\beta^\circ$ de $\mathbb{R}\mathcal{H}om^\circ(u_1^*(F_\bullet), u_2^!(G_\bullet))$ définit un élément
\[
\overline{\beta^\circ} \in \mathrm{Corr}(Y,G;\, X,F).
\]
\note{pour $\beta^\circ$ la page répète les mêmes arguments que pour $\alpha^\circ$ (les indices de $u$ ne sont pas sûrs) ; on attendrait $(u_2^*(G_\bullet), u_1^!(F_\bullet))$. L'argument se poursuit au-delà de ce lot.}

\end{document}
